% ======================================================================
% 07-discussion.tex
% ======================================================================
\section{Discussion: what AI did and did not do}
\label{sec:discussion}

\subsection{Separate the harness from the model}
\label{sec:harness-model}

Public discussions of AI-assisted programming usually compare models on a single
completion: given a prompt, which one writes the better function? That is not the
unit that mattered here. The unit was a \emph{turn}: read the repository, form a
hypothesis, edit, run the gate, record the result. Pi supplied the loop; DeepSeek
supplied the hypothesis. Most of the leverage came from the loop.

Three consequences follow. First, a model that is slightly worse at code
generation but equally good at following a procedure would have produced almost
the same project, because the procedure, not the completion, was the bottleneck.
Second, the value compounds with the repository: the more memory the agent can
read, the fewer turns are wasted re-learning. Third, the failure modes are
process failure modes, not syntax failure modes---a plausible hypothesis that
the gate rejects is the normal case, not an exception.

\subsection{Where the agent was useful}
\label{sec:agent-good}

\begin{itemize}
  \item \textbf{Repetition.} The five-seed $\tau$-ramp sweep exists because running
        it four more times was cheap. The 400-scenario certificate exists
        because re-scoring on a fresh seed was a command, not a project.
  \item \textbf{Parity and regression tests.} Comparing two implementations
        field by field, at several tolerances, on many scenarios, is tedious
        and exactly the kind of work an agent does without drift.
  \item \textbf{Documentation from artifacts.} The run log, the runtime table
        and the conclusions were rebuilt from stored JSON rather than memory,
        which is why the numbers in this article can be traced.
  \item \textbf{Searching the history.} Reconstructing when a claim changed
        across 69 commits is a task a human postpones and an agent performs.
  \item \textbf{Writing the guard.} After each failure, the agent wrote the
        smallest test that would have caught it. The guard is often four lines;
        the discipline is remembering to write it.
\end{itemize}

\subsection{Where the agent was wrong}
\label{sec:agent-bad}

The failure catalogue is also a catalogue of agent errors, and it is worth
naming them:
\begin{itemize}
  \item \textbf{Confident single-seed claims.} The $\tau$-ramp result was reported as
        an improvement before the paired sweep. The prose was plausible; the
        evidence was one seed.
  \item \textbf{Plausible structural fixes.} The FNO fix required proving that
        the old architecture \emph{could not} represent a front. A patch that
        merely trained longer would have been plausible and wrong.
  \item \textbf{Over-broad cleanups.} A ``simplification'' that removes the
        residual chunking is a regression on the 961\,MB host, even though it is
        correct on a GPU box. The memory file had to say so explicitly.
  \item \textbf{Metric drift.} Editing an evaluation script can silently change
        what a number means. The manifests and the artifact paths are the
        defence; a diff review alone is not.
  \item \textbf{Optimism about hardware.} The $5\times$ speedup and the A100
        assumption were both plan-level claims that survived until they were
        instrumented.
\end{itemize}

None of these is exotic. They are the ordinary failure modes of a competent,
tired collaborator. The response is therefore procedural rather than
technological: make the gate cheap, make the number traceable, and make the
negative result as easy to record as the positive one.

\subsection{The human protocol}
\label{sec:protocol}

The project converged on five rules, and they are the part most likely to
transfer:
\begin{enumerate}
  \item \textbf{Set the threshold before seeing the result.} A gate chosen after
        the run is a description, not a test.
  \item \textbf{Demand a failing test first.} A bug fix without a test that
        fails on the old code is a hypothesis, not a repair.
  \item \textbf{Ask what would falsify the claim.} If the answer is ``another
        seed'', run it; if the answer is ``nothing'', the claim is not
        empirical.
  \item \textbf{Keep the negative result.} The $\tau$-ramp, the A100 and the 3-well
        policy are as valuable as the successes, because they stop the project
        from repeating them.
  \item \textbf{Own the final judgement.} The agent can draft the conclusion;
        the human decides whether the evidence class supports it.
\end{enumerate}

\subsection{Threats to validity}
\label{sec:threats}

This is an experience report, not a controlled experiment. The sample is one
project, one human and one agent configuration; the models changed during the
work; there is no baseline project developed without AI; and the author is not
neutral about the outcome. The quantitative results are reproducible from the
repository, but the \emph{causal} claim---that the AI loop caused the
reusability---is an interpretation, supported by the commit history and the
failure catalogue, not by randomization.

Two further caveats are specific to the science. The 3D result is validated at
three presets but not against an independent 3D solver; the 2D certificate is
strong on reliability but shows that the reported CVaR tail does not survive a
fresh draw; and the 1D/FNO medians are five seeds, which is enough to see a 7$\times$
spread but not enough to estimate a tail. These are the same caveats the
repository records in \code{Conclusions.md}, and they are repeated here because
an article that hides them would reproduce the failure it is describing.

\subsection{What would make this a stronger study}
\label{sec:future}

A controlled comparison would need two teams, the same task and the same
deadline, one with and one without an agent, and a pre-registered measure of
reusability (e.g.\ time to reproduce a result on a fresh machine, number of
claims that survive an out-of-sample check). That study is worth doing. In its
absence, the best this article can do is make its own evidence trail complete
enough that a reader can disagree with the interpretation while trusting the
numbers.
